% The corpus instance: sizing a hundred-gigabyte open-text instance to study a knowledge store at scale.
\documentclass[11pt]{article}
\usepackage[T1]{fontenc}
\usepackage{lmodern}
\usepackage{microtype}
\usepackage[margin=1in]{geometry}
\usepackage{amsmath,amssymb,amsthm,mathtools}
\usepackage{algorithm}
\usepackage{algpseudocode}
\usepackage{booktabs}
\usepackage{array}
\usepackage[numbers,sort&compress]{natbib}
\usepackage[hidelinks]{hyperref}
\usepackage[capitalise,noabbrev]{cleveref}

\theoremstyle{plain}
\newtheorem{theorem}{Theorem}[section]
\newtheorem{lemma}[theorem]{Lemma}
\newtheorem{proposition}[theorem]{Proposition}
\newtheorem{corollary}[theorem]{Corollary}
\theoremstyle{definition}
\newtheorem{definition}[theorem]{Definition}
\theoremstyle{remark}
\newtheorem{remark}[theorem]{Remark}

\newcommand{\lean}[1]{\texttt{#1}}
\newcommand{\code}[1]{\texttt{#1}}
\newcommand{\Neff}{N_{\mathrm{eff}}}
\DeclareMathOperator{\Var}{Var}

\title{Sizing a Hundred-Gigabyte Instance of a Knowledge Store:\\
Information Never Hurts, Retrieval Crowds, and What a Gigabyte Costs}
\author{Alberto Serrano-Calva}
\date{\today}

\begin{document}
\maketitle

\begin{abstract}
A knowledge store that builds context from many sources should be able to show, at a scale where
its macro quantities are stable, what added context does for its answers and what each source
contributes. We design an instance of such a store that grows toward a hundred gigabytes of open,
continuously updated text, and we size it from measurements rather than estimates. Two pieces of
theory fix the design. First, conditioning never increases entropy, $I(X;Y,Z)\ge I(X;Y)$, so an
ideal reader can only gain from added context; the practical failure is finite retrieval, and we
prove the one-line lemma that under a fixed scorer and a fixed budget $k$ the top $k$ of a larger
corpus restricted to a smaller one lies inside the smaller one's top $k$, so recall is nonincreasing
in the corpus. Ranking must therefore select from structure that scales with the corpus: scopes,
dense cores, flow potentials, provenance per source. Second, the effective number of independent
signals among $N$ signals of mean correlation $\bar\rho$ is $N/(1+(N-1)\bar\rho)\le 1/\bar\rho$,
while across $D$ weakly correlated domains it grows as $D/\rho_w$; varied sources, and the bridges
between them, are what reach the regime in which macro behaviour is visible. On six sources
(Wikipedia, OpenAlex, arXiv, Hacker News, Stack Exchange, Project Gutenberg) we measure that prose
compresses 2.7 to 3.1 times at \code{zstd -3}, that a gigabyte of text is 1.1 to 1.7 million chunks of
1,200 characters, and that the local embedder yields 26 chunks a second on one stream and about 30
at saturation while raising every other caller's latency from 23 ms to seconds. A hundred gigabytes
is therefore 134 million vectors: 206 GB in float16 at 768 dimensions, 13 GB as one-bit codes, and
under a 92 GB disk cap only the code fits, with exact rescoring by re-embedding the candidates'
text at 1.5 s a query; exact float16 vectors fit up to 33 GB of text. The live store's row shapes
cost 1,084 bytes an occupancy and 567 bytes an edge and would exceed a terabyte; a columnar graph
is 31 GB. Embedding is the bottleneck at sixty days of the embedder for the last stage. We give a
three-stage protocol at 1, 10 and 100 GB with the measurements each must produce and what would
falsify each expectation, a source catalogue with every licence read, and the owner's decisions
with their prices. Nothing is ingested; the measurements are read-only.
\end{abstract}

\section{Introduction}
A companion study measured a single growing citation graph and found where its macro quantities
stop fluctuating: the spectral gap of a dense core concentrates once the core holds one to two
thousand nodes with a minimum degree near $\ln N$, and the Hodge split of the dated citation flow
settles one to two decades of $N$ earlier. A private store of a few dozen notes is below those
thresholds by orders of magnitude, so nothing measured on it about its macro regime is evidence.
The question here is how to build a store large and varied enough that the macro behaviour of the
knowledge system, not of one graph, can be studied: what context does for answers as the corpus
grows, where the system's constructs earn their keep, and what the system costs per gigabyte.

The direction the design follows is that adding context is good, connected context and disparate
trends both, and that the design's job is to make that true in practice. The contributions:
\begin{enumerate}
  \item two statements of theory made precise, with what follows for the design (\cref{sec:theory}):
    information never hurts an ideal reader and finite retrieval crowds (a lemma), and the
    effective number of independent signals is capped near $1/\bar\rho$ within a domain and grows
    with the number of weakly correlated domains;
  \item the design of a named instance, isolated from the default one, with one writer and every
    item labelled by its source as observed, and a batched bulk path through the store's ingestion
    language (\cref{sec:design});
  \item the storage and compute arithmetic from measurements: compression, chunk counts, the bytes
    per vector of five options with what each costs in retrieval quality, the bytes per row of the
    live graph stores, and the embedder's throughput and its effect on other callers
    (\cref{sec:arithmetic});
  \item a three-stage protocol with the measurements each stage must produce and the falsifier of
    each expectation, a source catalogue with every licence read, and the open decisions with
    their prices (\cref{sec:protocol,sec:catalogue}).
\end{enumerate}
Nothing is built and nothing is ingested. The samples were fetched into a scratch directory and
measured there; the live stores were read, not written.

\section{Background and related work}
\paragraph{Information and context.} The chain rule for mutual information and the nonnegativity
of conditional mutual information give $I(X;Y,Z)\ge I(X;Y)$, and conditioning reduces entropy
\citep[Theorems 2.5.2 and 2.6.5]{cover2006}. Retrieval-augmented readers condition a language model
on retrieved passages \citep{lewis2020rag}; that such readers use a long context unevenly, with
accuracy falling when the relevant passage sits in the middle, is measured by \citet{liu2024lost}.
The lemma of \cref{sec:crowding} is not about the reader's use of its window but about what reaches
the window: it is a statement about top-$k$ selection over a growing set.

\paragraph{Effective sample size.} The variance of a mean of $N$ equicorrelated variables is
$(1+(N-1)\bar\rho)/N$; the factor $1+(N-1)\bar\rho$ is Kish's design effect \citep{kish1965}, and
$N/(1+(N-1)\bar\rho)$ the effective sample size. The block form of \cref{sec:effective} follows by
summing the covariance matrix by blocks.

\paragraph{Vectors and codes.} Product quantization \citep{jegou2011pq}, binary codes with a
theoretical error bound \citep{gao2024rabitq}, graph indexes \citep{malkov2020hnsw} and
Matryoshka representations that keep a prefix of the dimensions useful \citep{kusupati2022matryoshka}
are the levers a vector plane has against its size. The embedder here is EmbeddingGemma-2
\citep{vera2025embeddinggemma}, 768 dimensions, stored in float16 \citep{ieee754} in LanceDB
\citep{lancedb}; the store's VEC-1 bounds what float16 can move a cosine by and certifies a top $k$
whose margin exceeds the bound. The recall of each code on the live corpus was measured on
2026-10-30 and is cited in \cref{sec:vectors}.

\paragraph{Scale thresholds and attribution.} The thresholds SCL-1 to SCL-4 (the relaxation time,
the eigengap at the cut, the concentration of spectra once the minimum degree exceeds $C\log N$
\citep{oliveira2009concentration,chung2011spectra}, the power calculation), the $k$-core
\citep{seidman1983cores}, HodgeRank \citep{jiang2011hodgerank}, and the Shapley attribution
\citep{shapley1953value,castro2009polynomial} with its leave-one-out contrast are the companion
notes' and are used, not restated. Bridges between communities \citep{granovetter1973weak}
and community structure \citep{newman2006modularity,fortunato2010community} are the objects the
cross-source measurements count. Link prediction on a store's own history
\citep{libennowell2007link} is the ledger's first generator.

\paragraph{Compression.} Zstandard \citep{rfc8878} is the raw store's codec; its levels are the one
lever on the raw bytes.

\section{Theory}\label{sec:theory}

\subsection{Information never hurts; finite retrieval crowds}\label{sec:crowding}
Let $X$ be what a question asks for, $Y$ the context a reader has and $Z$ context added.

\begin{proposition}[conditioning never increases entropy]\label{prop:info}
$I(X;Y,Z) = I(X;Y) + I(X;Z\mid Y) \ge I(X;Y)$, equivalently $H(X\mid Y,Z)\le H(X\mid Y)$. For any
bounded loss, the Bayes risk of a decision on $(Y,Z)$ is at most that of a decision on $Y$.
\end{proposition}
\begin{proof}
The equality is the chain rule; the inequality is $I(X;Z\mid Y)\ge 0$ \citep[2.5.2, 2.6.3]{cover2006}.
The decision on $Y$ alone is among the decisions available on $(Y,Z)$, so the minimum over the
larger set is no larger.
\end{proof}

So added context, connected or disparate, cannot hurt an ideal reader. If it hurts, the reader is
not ideal, and the one way this system's reader is not ideal is that it reads through a finite
retrieval: a scorer ranks the corpus and the reader sees the $k$ best.

\begin{lemma}[crowding]\label{lem:crowding}
Let a scorer induce a total order on items (ties broken consistently) and $R_k(C)$ be the $k$
best items of a finite set $C$. For $C\subseteq C'$,
\[
  R_k(C')\cap C \subseteq R_k(C).
\]
Hence for a relevant item $r\in C$, $\mathbf 1[r\in R_k(C')] \le \mathbf 1[r\in R_k(C)]$: under a
fixed scorer and a fixed budget, recall at $k$ is nonincreasing as the corpus grows, whatever is
added.
\end{lemma}
\begin{proof}
An item of $C$ among the $k$ best of $C'$ is beaten by fewer than $k$ items of $C'$, so by fewer than
$k$ items of $C\subseteq C'$, so it is among the $k$ best of $C$.
\end{proof}

The reader receives $I(X;R_k(C'))$, not $I(X;C')$; the first can fall while the second rises, and
the difference is the crowding: the relevant piece is pushed out of the window by pieces that score
above it and inform less. The remedy is not less context but a budget that scales with the corpus:
the ranking must select from a structure, not from one flat list. If a scope $\sigma$ contains $r$
and is sized by the question rather than by the corpus, the recall inside $\sigma$ is that of
$|\sigma|$, not of $|C'|$, and \cref{lem:crowding} applied to $\sigma\subseteq C'$ says the scoped
read is never worse than the unscoped one for an $r\in\sigma$. The store has four such structures:
scopes (\code{over} a source, \code{within} a region, a congruence on kernels that pushes to the
leaves), dense cores (the $k$-core of a community, where the companion study found the gap
concentrates and the whole component's never does), flows (the HodgeRank potential of a dated
relation, an order that is not the scorer's and does not crowd with the corpus), and provenance per
source (the origin every item carries, so a reader can ask of one source and the ledger can
attribute to it). The empirical questions are the rate at which unscoped recall falls with $N$ on a
real corpus and whether the scoped recall is flat; they are stage 1's battery (\cref{sec:protocol}).

\subsection{The effective number of independent signals}\label{sec:effective}

\begin{proposition}[Kish]\label{prop:kish}
Let $N$ signals have unit variance and mean pairwise correlation $\bar\rho\in[0,1]$. The variance of
their mean is $(1+(N-1)\bar\rho)/N$, so the effective number of independent signals is
\[
  \Neff = \frac{N}{1+(N-1)\bar\rho}\le\frac1{\bar\rho},\qquad \Neff\to\frac1{\bar\rho}\ \text{as}\ N\to\infty.
\]
\end{proposition}
\begin{proof}
The sum of the covariance matrix's $N^2$ entries is $N + N(N-1)\bar\rho$; divide by $N^2$
\citep{kish1965}.
\end{proof}

\begin{proposition}[domains]\label{prop:domains}
With $D$ domains of $m$ signals each, correlation $\rho_w$ within a domain and $\rho_b$ between
domains,
\[
  \Var(\text{mean}) = \frac{1+(m-1)\rho_w + m(D-1)\rho_b}{Dm},\qquad
  \Neff \to \frac{D}{\rho_w + (D-1)\rho_b}\ \text{as}\ m\to\infty,
\]
which is $D/\rho_w$ when the domains are uncorrelated and $1/\rho_w$ when they are as correlated as
their members.
\end{proposition}
\begin{proof}
Sum the covariance matrix by blocks: $Dm$ diagonal entries, $Dm(m-1)$ within-block off-diagonal
entries at $\rho_w$, $Dm\cdot m(D-1)$ between-block entries at $\rho_b$; divide by $(Dm)^2$ and let
$m\to\infty$.
\end{proof}

Read through the companion note's thresholds: SCL-4's $K$, the number of independent resolved
units an effect $\Delta$ needs, $K\ge((z_{1-\alpha/2}+z_{1-\beta})s/\Delta)^2$, is bounded by $\Neff$.
A corpus of one domain cannot push $K$ past about $1/\bar\rho$ however large it grows; the macro
behaviour of the system, a difference between arms at the bar, becomes visible only when $D/\rho_w$
clears $K$. Varied sources are what reach that regime, and the connections between domains, the
bridging edges between communities of different sources and the rises they share, are themselves
the quantities to measure: $\rho_b$ is read off them. A signal here is a community's activity
series at the instance's cuts; $\rho_w$ and $\rho_b$ are its correlations within and across
sources. The hypothesis the instance tests is that $\rho_b$ on these six sources is well below
$\rho_w$.

\section{Concepts and intuition}
Three pictures carry the design. The first is a window: the reader's context is a window of $k$
passages over a corpus that grows under it; \cref{lem:crowding} says the window can only lose a
given passage as the corpus grows, and the store's constructs are ways of placing the window over
a part of the corpus whose size the question sets. The second is a sum of correlated coins: a
thousand signals from one domain are worth about $1/\bar\rho$ independent ones, a few dozen at
typical correlations, and only another domain adds another few dozen; so a store that wants to see
macro behaviour needs breadth before depth, and the bridges between its domains are where the
breadth is tested. The third is a budget line: a hundred gigabytes of text is a hundred and thirty
million vectors, and the disk under them is the same disk the private store runs on, so every byte
a vector keeps is a byte the raw text, the graph or the headroom loses.

\section{Design}\label{sec:design}

\subsection{A named instance}
The instance's id is \code{corpus}. Under the store's instance identity (DI-2) everything it holds
derives from that id and is disjoint from the default instance by construction: its root is
\code{instances/corpus/}, its stores, spool, lock and levers under it, its secret names and service
labels prefixed by it, its exhaust lane inside its root. It declares no vault, no transcripts and
no code agent, so it reads none of the default instance's sources; the loader refuses an overlay
that could reach anything another instance writes, and a property test over generated ids checks
the invariant. It has one writer, its supervisor; every surface writes a request its daemon
enqueues. Every item lands through a script of the ingestion language under the label its source's
rule gives, and that rule is the same for all six sources: observed. What distinguishes an item is
its origin, the source name and the item's path, which the value carries and no expression can
change; so the store can say of every chunk which source it came from, scoped reads can restrict
to a source, and the attribution can ablate one. A version of an item (a revision, an arXiv
\code{v2}, an edited answer) is a new version on the identity's chain, never a new identity.

\subsection{The disk}
The instance shares one disk with the default instance. 117 GB are free and 25 GB must stay free,
so the instance's whole footprint is capped at 92 GB, below the 100 GB dataset budget. Every figure
below is read against 92.

\subsection{The bulk path}\label{sec:bulk}
The daemon's job rate, read from its queue, is 720 an hour, 12 a minute: the tick's cadence, not
the writer's ceiling. A path that enqueued one item per job would land $10^7$ items in 580 days; the
bulk path makes a job a batch. Each source is a line of the ingestion language's table: a structured
source with a location under the instance's root, a format, an item rule, a label rule (observed),
a refresh (an edge process that fetches a dump part or drains a stream into the source's directory
and writes batches of items as files the daemon trusts) and a cursor (the field whose newest value
is recorded once a plan that read the source has landed). The script is one per source and the same
shape for all six (\cref{alg:sync}); the one addition the language needs is a \code{batch} argument
to its \code{items} reading function, so a plan is bounded in gas and lands as one command of the
keyed model. By the ingestion laws, batches on distinct identities land in any order and converge
under re-run after any interruption (ING-9, ING-11), and a re-fetched item whose text the plane
holds embeds nothing (ING-5), which makes a daily stream drain idempotent. A batch of 2,000 items
is about 2,700 chunks, 100 s of the embedder, one transaction; $10^7$ items are 5,000 jobs, seven
hours at the tick's rate, and the embedder, not the queue, is what the jobs wait on.

\begin{algorithm}
\caption{One batch of a corpus source, observed, under its own origin.}\label{alg:sync}
\begin{algorithmic}[1]
\Require a source name and a batch size $b$
\Ensure the next $b$ items past the source's cursor landed, the cursor advanced once the plan lands
\For{each item $it$ in $\mathrm{items}(\mathrm{source}(\text{``corpus/''}+name),\ \mathrm{batch}{:}\ b)$}
  \State store $it.\mathrm{content}$ in ``observed'' at ``corpus/''$+name+$``/''$+it.\mathrm{id}$
\EndFor
\end{algorithmic}
\end{algorithm}

\section{The measured arithmetic}\label{sec:arithmetic}

\subsection{Method}
Samples of the six sources were fetched into a scratch directory (88 MB of data and a 21 MB
catalogue): 24.1 MB of bz2 streams cut on stream boundaries from part 1 of the current Wikipedia
dump (2,100 pages, 1,602 articles after redirects); 800 questions and 800 answers of one Stack
Exchange site through its API; 2,000 OpenAlex works of 2024 with abstracts through its API, the
abstract rebuilt from its inverted index; one OAI-PMH page of 1,300 arXiv \code{cs} records;
3,000 comments and 1,000 stories of Hacker News through the Algolia API; twenty Project Gutenberg
texts. Each was reduced to the form that would land (prose with markup removed, a title prepended)
and measured: UTF-8 bytes, \code{zstd} at levels 3, 9 and 19 on the concatenation, and the store's
own chunker (1,200 characters, 150 of overlap, block-aware) for the chunk count. The embedder was
timed against the stack's service (EmbeddingGemma-2 Q8\_0, four slots, context 8,192, all layers on
the GPU) with the chat model checked idle through its slots endpoint before every configuration,
six configurations of 512 chunks each, 3,000 chunks in all drawn interleaved from five sources
(973 characters, 240 tokens each by the server's count); a single-query probe was timed idle and
once beside each configuration. The live stores were read, not written, for their bytes per row.

\subsection{Compression and chunks}
\begin{table}[t]
\centering
\caption{Measured ratios per source. Compression ratios are uncompressed over compressed bytes of
the landing-form text. Chunks are the store's chunker at 1,200 characters with 150 of overlap.}\label{tab:ratios}
\footnotesize
\begin{tabular}{lrrrrrr}
\toprule
source & items & item (B) & zstd $-3$ & zstd $-19$ & chunks/MB & chunk (chars) \\
\midrule
Wikipedia, prose & 1,602 & 22,531 & 2.77 & 3.59 & 1,217 & 964 \\
Stack Exchange, posts & 1,600 & 1,442 & 2.95 & 3.64 & 1,313 & 830 \\
OpenAlex, title and abstract & 2,000 & 1,588 & 3.08 & 3.94 & 1,578 & 704 \\
arXiv, title and abstract & 1,300 & 1,514 & 2.93 & 3.58 & 1,728 & 649 \\
Hacker News, one comment an item & 2,000 & 251 & 2.65 & 3.00 & 4,212 & 240 \\
Project Gutenberg, books & 20 & 932,137 & 2.83 & 3.67 & 1,129 & 1,018 \\
\bottomrule
\end{tabular}
\end{table}

\Cref{tab:ratios} gives the ratios. Prose compresses 2.7 to 3.1 times at \code{zstd -3} and 3.6 to
3.9 at $-19$; the level buys a quarter more at about ten times the compressor's time. A short item
costs a whole chunk however short, so comments taken one at a time make 4,200 chunks per MB; the
design takes a Hacker News thread as the item, each comment a block, and sizes it at the Stack
Exchange rate, an assumption stage 1 measures. Other counts from the samples: a Wikipedia article in
the slice links to 205 distinct pages on average (median 133); an OpenAlex work of 2024 with an
abstract references 135 works on average (median 70), 93\% at least one.

\begin{table}[t]
\centering
\caption{The 100 GB mix and what it becomes. Text is UTF-8 in the landing form; chunks and
compressed sizes follow \cref{tab:ratios}. Wikipedia is every article's prose and a 5 GB history
subset; OpenAlex 15 M works with abstracts in selected fields; Stack Exchange five sites with
comments and edits; Hacker News every item as threads, chunked at the Stack Exchange rate; Project
Gutenberg 10,000 English texts; arXiv every record's title and abstract with its versions.}\label{tab:mix}
\small
\begin{tabular}{lrrrr}
\toprule
source & text (GB) & chunks (M) & zstd $-3$ (GB) & zstd $-19$ (GB) \\
\midrule
Wikipedia & 35.0 & 42.6 & 12.6 & 9.7 \\
OpenAlex & 24.0 & 37.9 & 7.8 & 6.1 \\
Stack Exchange & 15.0 & 19.7 & 5.1 & 4.1 \\
Hacker News & 12.0 & 15.8 & 4.5 & 4.0 \\
Project Gutenberg & 9.5 & 10.7 & 3.4 & 2.6 \\
arXiv & 4.5 & 7.1 & 1.5 & 1.1 \\
\midrule
total & 100.0 & 134.4 & 34.9 & 27.7 \\
\bottomrule
\end{tabular}
\end{table}

\Cref{tab:mix} spends the budget where the instruments need it, with dated graphs and version
chains in every source but one. A hundred gigabytes of text is 134 million chunks and 35 GB of raw
store.

\subsection{Vectors}\label{sec:vectors}
\begin{table}[t]
\centering
\caption{Bytes per vector of five options, their total at each stage, and what each keeps of the
exact top ten on the live corpus (recall at 10 against exact float32, measured 2026-10-30 on the
titles, owner and self query families).}\label{tab:vectors}
\footnotesize
\begin{tabular}{lrrrr>{\raggedright\arraybackslash}p{5.6cm}}
\toprule
option & B/vector & 1 GB & 10 GB & 100 GB & keeps of the exact top ten \\
\midrule
float16, 768 & 1,536 & 2.1 GB & 20.5 GB & 206 GB & 0.999 to 1.000; VEC-1 certifies the margin \\
int8, 768 & 768 & 1.0 & 10.3 & 103 & 0.989 to 0.993; an index beside the float column only \\
float16, 256 (Matryoshka) & 512 & 0.7 & 6.8 & 69 & 0.46 to 0.76; 0.78 to 0.96 rescored from full vectors \\
int8, 256 (Matryoshka) & 256 & 0.3 & 3.4 & 34 & not measured; at most the row above \\
one-bit code, 768 bits & 96 & 0.1 & 1.3 & 13 & sign bits 0.17 to 0.60; 1-bit RaBitQ refined 0.86 to 0.97 \\
\bottomrule
\end{tabular}
\end{table}

Each chunk is one vector of 768 dimensions. \Cref{tab:vectors} gives the bytes and the measured
recall of each option. VEC-1 bounds the cosine error of float16 storage at $4.884\times10^{-4}$ for
$d=768$ and certifies a top $k$ whose margin exceeds twice the sum of the storage, query and kernel
bounds; the live store certified 54\% of its queries. With the raw store and the compact graph of
\cref{sec:graph} beside them, the footprints at 100 GB are in \cref{tab:fits}: only the one-bit code
fits under 92 GB. Per gigabyte of text the footprint is 0.35 GB raw, 0.31 GB graph, 0.06 GB chunk ids
and 0.01 GB index plus the vectors, so exact float16 vectors admit 33 GB of text, int8 52, Matryoshka
float16 65, Matryoshka int8 85, and the one-bit code 106.

\begin{table}[t]
\centering
\caption{Footprints at the 100 GB stage, chunk grain, raw at \code{zstd -3}, against the 92 GB cap.}\label{tab:fits}
\small
\begin{tabular}{lrrrrrrl}
\toprule
configuration & raw & vectors & graph & ids & index & total (GB) & fits \\
\midrule
float16 768 & 35 & 206 & 31 & 6 & 1 & 280 & no \\
int8 768 & 35 & 103 & 31 & 6 & 1 & 177 & no \\
float16 256 & 35 & 69 & 31 & 6 & 1 & 142 & no \\
int8 256 & 35 & 34 & 31 & 6 & 1 & 108 & no \\
one-bit code & 35 & 13 & 31 & 6 & 1 & 86 & yes \\
float16 768, item grain only (36 M items) & 35 & 56 & 31 & 0 & 0 & 122 & no \\
\bottomrule
\end{tabular}
\end{table}

What makes the one-bit code usable at its measured unrescored recall is that the raw text is on
disk: the top $m$ candidates by code can be re-embedded at query time and rescored exactly, at
$m/26$ seconds (1.5 s for $m=40$), with no full-precision column kept; the VEC-1 certificate then
applies to the rescored list exactly as to a float16 store. This is the design's default for the
100 GB stage, and the recall it keeps on these sources is stage 2's first measurement, against
exact answers that still fit there.

\subsection{The graph and its indexes}\label{sec:graph}
The live stores' bytes per row, read read-only: a curated Lance row 2,215 B (the float16 vector
1,536 and 679 for the text column and ids), a plane occupancy 1,084 B with its three indexes, a
reference edge 567 B with its five, a membership 385 B. At the corpus's counts, 36 million items and
about 1.8 billion dated in-corpus edges (Wikipedia 7.3 M articles at about 150 article links each,
OpenAlex 15 M works at 40 in-corpus references, Hacker News 5 M threads at 9 replies, Stack
Exchange 6 M posts at 2.5, arXiv 3 M at 1.5 versions), the live shapes do not scale
(\cref{tab:graph}). The compact column is the design: the chunk text is not duplicated into the
vector table (the raw store holds it; a chunk is an item digest with a start and an end),
occupancies are at the item grain with integer ids and one index, edges are columnar per source
and per cut and are loaded as sparse matrices for the spectra, and only the quotient graph between
communities and the derived backbone land in the store's relation tables. The 120 B occupancy is a
target; stage 1 measures it.

\begin{table}[t]
\centering
\caption{The graph and its indexes at 100 GB: the live row shapes against the compact design.
Occupancies are counted at the chunk grain today and at the item grain in the compact design;
the last two columns are gigabytes.}\label{tab:graph}
\footnotesize
\begin{tabular}{lr>{\raggedright\arraybackslash}p{4.8cm}rr}
\toprule
component & count & unit, today / compact & today & compact \\
\midrule
chunk rows beside the vector & 134 M & 679 B with the text / 48 B as ids and offsets & 91 & 6.5 \\
plane occupancies & 134 M / 36 M & 1,084 B at the chunk grain / 120 B at the item grain & 146 & 4 \\
edges & 1.75 G & 567 B as indexed rows / 12 B as columnar triples & 994 & 21 \\
versions & 47 M & 120 B & 6 & 6 \\
IVF partition index & 134 M & 8 B & & 1 \\
\bottomrule
\end{tabular}
\end{table}

\subsection{Throughput and compute}
\begin{table}[t]
\centering
\caption{Embedding throughput against the local embedder, 3,000 chunks of 973 characters and 240
tokens, the chat model idle. The probe is a single short query timed beside the running
configuration; idle it took 23 ms.}\label{tab:embed}
\small
\begin{tabular}{rrrrrr}
\toprule
batch & streams & chunks/s & tokens/s & batch latency p50 (s) & probe beside it (s) \\
\midrule
8 & 1 & 21.4 & 5,298 & 0.34 & 0.60 \\
32 & 1 & 23.4 & 5,940 & 1.20 & 1.23 \\
64 & 1 & 26.1 & 6,245 & 2.38 & 1.84 \\
32 & 2 & 28.9 & 6,753 & 2.15 & 1.65 \\
32 & 4 & 30.4 & 6,795 & 3.99 & 4.32 \\
64 & 4 & 29.0 & 6,750 & 8.06 & 8.45 \\
\bottomrule
\end{tabular}
\end{table}

\Cref{tab:embed} gives the embedder's throughput. It saturates near 30 chunks/s (about 6,800
tokens/s); batch 64 on one stream gives 26, a second stream buys 11\% and four buy 16\%, and every
stream adds its queue to the latency another caller sees, which is the default instance's daemon
and the reader at the chat model. The design embeds on one stream at batch 64 and yields: a bulk
batch is not started while the embedder has another principal's request in flight. The compute per
stage (\cref{tab:compute}) is then 14 hours at 1 GB, six days at 10 GB and sixty days at 100 GB, the
same for every vector option, because a code is derived from the full vector; the embedder sleeps
after a minute idle and shares the GPU with a 27B chat model that has been measured to run out of
GPU memory under load, so the bulk path runs when the chat model is asleep and the sixty days are
calendar weeks at the machine's idle hours. Compression and chunking run at hundreds and tens of
MB/s and are minutes per stage. The daily drains (3.2 English Wikipedia edits a second measured on
the event stream, 10,000 Hacker News items a day, a few hundred Stack Exchange posts a day, about
2,000 arXiv records a day) are under 50,000 chunks a day, half an hour of the embedder.

\begin{table}[t]
\centering
\caption{Compute per stage at the measured single-stream rate.}\label{tab:compute}
\small
\begin{tabular}{lrrrr}
\toprule
stage & chunks (M) & embedder at 26/s & at 30/s & jobs of 2,000 items (12 a minute) \\
\midrule
1 GB & 1.3 & 14 h & 12 h & 181 (0.3 h) \\
10 GB & 13.4 & 143 h, 6 days & 5 days & 1,814 (2.5 h) \\
100 GB & 134 & 1,430 h, 60 days & 52 days & 18,140 (25 h) \\
\bottomrule
\end{tabular}
\end{table}

\section{The staged protocol}\label{sec:protocol}
Every stage is a dated extension of the one before, and every stage produces the same families of
measurement against $N$, so the curves are read across three decades. The method for the macro
quantities is the companion study's: yearly cuts, the smallest twenty eigenvalues of the normalized
Laplacian by shift-invert ARPACK with a minimum-degree ordering, node and edge bootstrap widths,
the 10\% crossing, the $k$-cores, the Hodge split of the dated flow by LSMR on the clique complex.
It ran at $N=33{,}000$ in minutes; the corpus's graphs are measured per source and per community at
that scale, and the whole is read on the quotient graph of the communities (DYN-5: exact on block
fields, the gap never lowered), never on 36 million nodes at once.

\paragraph{Stage 1: 1 GB} (170 MB a source, 1.3 M chunks, 14 h of embedding; everything exact,
float16 chunk vectors with every code beside them for the bench). The measurements: the macro
quantities against $N$ per source ($\gamma$, the first four eigengaps, $d_{\min}$, the $k$-core
crossings, the Hodge split at yearly cuts with bootstrap widths), six curves for six relation
kinds, falsified by a source whose 6-core or 10-core gap width never falls under 10\% by $N=10^5$;
the Hodge split's earlier settling per source, falsified by a harmonic fraction still moving where
$\gamma$ has settled; cross-source bridges (explicit: Wikipedia citations of arXiv and DOI works,
Hacker News and Stack Exchange links to Wikipedia and arXiv, the OpenAlex–arXiv twins; semantic:
chunk pairs across sources above $\sigma$) against a degree-preserving rewiring null, and shared
rises ($\rho_w$, $\rho_b$ at lags of 0 to 8 weeks, $\Neff$ by \cref{prop:domains}), falsified by
bridges at the null or $\rho_b$ within the band of $\rho_w$; the retrieval battery (500 held-out
titles a source, the labelled twins across sources, the BEIR pair of the embedder bench as a fixed
control \citep{thakur2021beir}; recall at 10 unscoped and \code{over} a source), falsified by an
unscoped recall flat in $N$ across the stages or a scoped recall falling with the source's own $N$;
attribution by the ledger's generators over the six sources with the ablation arms and the
Shapley reading, the OpenAlex–arXiv pair as the redundancy witness (leave-one-out near zero for
each, Shapley near a half each), falsified by both twins necessary at the bar; cost and latency
against $N$ at the profile's method, the occupancy bytes per item against the 120 B target, and
the Python share of each sentence, falsified by an occupancy above 300 B or a share above 20\%.

\paragraph{Stage 2: 10 GB} (13.4 M chunks, six days, 29 GB on disk; the last stage at which exact
float16 chunk vectors fit beside everything else, so it decides the vector plane). Everything of
stage 1 again, and: the codes against exact answers (int8, Matryoshka-256 in float16 and int8, the
one-bit code, each unrescored and rescored by re-embedding the top $m\in\{20,40,80\}$, against the
exact top ten on stage 1's battery, with the VEC-1 certified share and the latency of each),
falsified for the default by a one-bit code rescored at $m\le80$ below 0.95 on the titles and twins
families; an index (flat float16 at 13 M rows, about 1.8 s by linear scaling of the 19.7 ms measured
at 147 k, against IVF at nprobes 32 and 128), falsified by IVF at 128 below 0.95 on the twins; the
thresholds on the quotient graph of the communities and the relaxation time that spaces the
ledger's rounds, with DYN-5's bound as the control; SCL-4's $K$ from the pilot's spread against
the rounds a year of weekly cuts affords; and the compiled-rewrite verdict at 13 M rows.

\paragraph{Stage 3: 100 GB} (134 M chunks, sixty days, 86 GB on disk; the configuration stage 2
chose, by default one-bit codes with re-embedding, raw at \code{zstd -3} with $-19$ in reserve, the
compact graph, the full catalogue and the daily drains). Everything again across three decades of
$N$, and: where each quantity's width crossed 10\%, in $N$ and in $d_{\min}/\ln N$, against the
companion study's 1,000 to 2,000 and 0.9 to 1.3; $\Neff$ over all six sources and whether $K$ is
reached, falsified by an $\Neff$ no larger at 100 GB than at 10 (which \cref{prop:domains}
predicts for more text in the same domains, and only a seventh domain would lift); the crowding
curve over three decades and the $N$ at which unscoped recall on the twins falls under one half;
the Shapley and leave-one-out readings of the six sources with their disagreements named as
redundancy; and the daily cost (chunks a day, embedder minutes a day, disk growth a week against
the cap), falsified as a plan by growth past 1 GB a week.

In one line each: stage 1 says whether the sources connect and whether crowding is real at a
million chunks; stage 2 says what a code costs against exact answers and settles the vector plane;
stage 3 says whether the macro quantities, the effective count and the attribution behave across
three decades as \cref{lem:crowding}, \cref{prop:domains} and the companion study predict.

\section{The source catalogue}\label{sec:catalogue}
Every licence and term was read on 2026-11-07.

\begin{description}
  \item[Wikipedia] (English, current article text and a subset of revision histories). Text under
    CC BY-SA 4.0 and the GFDL \citep{wikimedia2026legal}; the Foundation's terms of use take
    precedence; automated clients must send a descriptive User-Agent with contact information or
    are refused. 7,260,766 articles; the current-articles multistream dump is 27.0 GB bz2 and over
    105 GB of XML, about 30 to 40 GB of prose by the measured ratio; the full history is multiple
    terabytes uncompressed in 7z parts. Dump parts fetched by byte range through the multistream
    index; the live edit stream over server-sent events, resumable by event id
    \citep{wikimedia2026eventstreams}, 3.2 English edits a second measured. It exercises the
    hyperlink graph, version chains with dates on every page, a velocity with many generations,
    and bridges to arXiv and DOI works through citation templates.
  \item[OpenAlex] (works with abstracts, selected fields). CC0 1.0 \citep{openalex2026snapshot,
    priem2022openalex}; the snapshot is free with no account; the API is 10 requests a second
    without a key, cursor paging needs a free key. 464,765,485 works in the 2026-10-14 manifest,
    675.3 GB gzipped in 1,898 parts; 168.7 M with abstracts, of which Mathematics 2.3 M, Computer
    Science 8.4 M, Physics and Astronomy 20.6 M. Quarterly snapshot streamed through a field filter
    and discarded; the API for the daily delta. It exercises the dated citation graph the companion
    study measured, topics as a partition, and the OpenAlex–arXiv twin as the redundancy witness.
  \item[arXiv] (metadata with abstracts). Descriptive metadata including abstracts under CC0 1.0
    \citep{arxiv2026tou}; e-prints are not and are not fetched; at most one request every three
    seconds. More than three million articles, 4.5 GB of text. OAI-PMH with versions and a licence
    per record, incremental by date, about 2,000 records a day. It exercises version chains with
    dates and no citations, and a nightly feed the loop registers from round by round.
  \item[Hacker News] The official API has no rate limit and no stated terms on the content
    \citep{hn2026api}; the search API allows 10,000 requests an hour; the site's terms grant the
    content's licence to its operator and no third party and forbid scraping the site, which the
    API is not; the content carries no open licence. 50,408,507 items; 9,300 comments and 850
    stories a day measured; about 12 GB of text as threads. It exercises a reply tree with
    minutes of latency, links to Wikipedia and arXiv as explicit bridges, and the thread as an item
    whose comments are blocks. Its inclusion is an owner's decision.
  \item[Stack Exchange] (math, physics, stats, cs, mathoverflow). User content under CC BY-SA (2.5,
    3.0, 4.0 by date) with attribution to the question and the author; the 2024-04-02 archive is
    CC BY-SA 4.0 with no login \citep{stackexchange2024dump}; later official dumps are downloaded
    from a logged-in profile under terms that admit personal use and projects that do not train a
    large language model; the API is 300 requests a day anonymous with a backoff field to honour.
    The five sites hold 2.4 M questions, 3.1 M answers and 13.6 M comments; about 15 GB of text with
    edit history. It exercises answers and comments as a dated relation, edits as version chains,
    and the question-and-answer form the battery's queries come in. The dump path is an owner's
    decision.
  \item[Project Gutenberg] (English texts). Public domain in the United States
    \citep{gutenberg2026policy}; the trademark may be kept or stripped for non-commercial use;
    copyrighted items are excluded by the catalogue's rights field; bulk access through the
    harvest URL with a two-second wait or a mirror, not the main site. 78,461 texts, 63,115 in
    English; the sample's books 0.93 MB each, so 10,000 texts are about 9.5 GB. It exercises long
    documents with no dated relation: the null for the velocity instruments and the control for
    the cores.
\end{description}
Nothing is dropped outright; two sources carry conditions the owner decides on, and both default
to the path that needs no account.

\section{The owner's decisions, priced}
\begin{description}
  \item[D-1, the vector plane at 100 GB.] Default: one-bit codes at the chunk grain with exact
    rescoring by re-embedding the top 40, no full column, 86 GB on disk; price: 1.5 s a query
    against 20 ms today, and a recall to be measured at stage 2. Alternatives: stop at 33 GB of
    text with exact float16 vectors (every answer certifiable, a third of the catalogue, six domains
    still); an external disk (exact at 100 GB is 280 GB); int8 Matryoshka-256 up to 85 GB of text
    (recall unmeasured).
  \item[D-2, Hacker News.] Default: included through the official API, held locally, never
    republished; price: the one source without an open licence. Alternative: drop it, returning
    12 GB and 16 M chunks to the others and losing the reply tree, the minutes-scale velocity and
    one of six domains.
  \item[D-3, Stack Exchange's dumps.] Default: the 2024-04-02 archive plus the anonymous API;
    price: the backfill since arrives at 300 requests a day. Alternatives: the owner's account and
    the current official dumps under their terms; a registered API key.
  \item[D-4, the embedder's sixty days.] Default: this machine, one stream, at idle hours; the
    vectors are keyed by content and embedder, so a pause loses nothing. Alternatives: a second
    machine with the same model; the item grain only (36 M vectors, 16 days, no chunk-level
    retrieval for long items); 2,400-character chunks (half the count, a different retrieval unit).
  \item[D-5, OpenAlex access.] Default: the quarterly snapshot streamed through a field filter, no
    account, 675 GB of transfer a quarter. Alternative: a free key for cursor paging and the daily
    delta.
  \item[D-6, the Wikipedia history subset.] Default: 5 GB of text, the revisions of the pages the
    other sources bridge to, sampled monthly. Alternatives: every revision of a thousand pages;
    none.
  \item[D-7, the raw compression level.] Default: \code{zstd -3}, 35 GB. Alternative: $-19$ from
    the start, 7 GB saved at ten times the compressor's time, minutes not hours.
\end{description}

\section{Limitations}
One sample per source and one run per embedder configuration: the ratios are good to one
significant figure and the throughput carries the machine's ordinary variance. The Wikipedia slice
is from the dump's first part, the oldest page ids, whose articles are long; its item size overstates
the encyclopedia's, its chunk rate per megabyte does not depend on it. The wikitext reduction is a
regular-expression pass, not a parser, so the prose ratio is slightly high. The Hacker News thread
rate and the compact occupancy's 120 bytes are assumptions the first stage measures. The recall of
a one-bit code rescored by re-embedding has not been measured on these sources; the default for
the last stage rests on the measured recall of the RaBitQ index refined from full vectors, which is
the same operation with a different source of the exact vectors. The licences were read on one day
and change; the catalogue records the date. The edge counts are means from the samples applied to
the sources' own totals, and the Wikipedia link count includes links the article namespace filter
would drop. Nothing here was ingested, so nothing here is a result about the store's behaviour at
scale; that is what the protocol is for.

\section{Future work}
The first fetch, licensed at the merge and following the ingestion language's merge, lands the
1 GB stage and its battery. The ingestion language gains the \code{batch} argument and the
instance its configuration section; six refresh processes land under the edge zone. The 10 GB stage
decides the vector plane by measurement, and if the one-bit code fails its bar the owner's
alternatives of D-1 are priced here already. A seventh domain, if $\Neff$ saturates as
\cref{prop:domains} predicts, is the way to lift it; candidates are archival economic series with
vintages and the commit histories of public repositories, both named in the companion note. The
lag structure of the shared rises, whether a rise in one source leads another's, is the HodgeRank
potential of a lead-lag flow and is the natural next reading once the series exist.

\section{Conclusion}
Added context cannot hurt an ideal reader and does hurt a finite retrieval, by a lemma whose remedy
is a ranking that selects from structure sized by the question; the effective number of independent
signals is capped within a domain and grows with the number of weakly correlated domains, so breadth
across sources, and the bridges between them, is what makes macro behaviour visible. A hundred
gigabytes of open text from six such sources is 134 million chunks, 35 GB compressed, and sixty days
of this machine's embedder; its vectors fit under the disk's cap only as one-bit codes rescored
exactly from the raw text, or the corpus stops at a third of the budget with exact vectors. The
three-stage protocol measures the thresholds, the bridges, the crowding, the codes and the
attribution against $N$ across three decades, with the falsifier of each expectation named before
the first byte is ingested.

\bibliographystyle{plainnat}
\bibliography{refs}
\end{document}
